\section{Seven Cross-Period Tests: What Actually Continues?}

The preceding history is broad by necessity, but breadth can conceal a new
error. Once biblical narratives, apologetic claims, initiatory prayers,
offices, manuscript books, Roman editions, codes, current law, clinical
classification, and safeguarding sources occupy one volume, their mere
proximity can look like a single tradition moving intact through time. The
opposite error is just as easy: every change of word, book, or authority can
be treated as rupture.

A responsible synthesis must name the object whose continuity or change is
being claimed. Seven tests do that work: recipient, act, minister, book,
authority, evidence, and care. They do not form a ritual sequence and must not
be used as an intake checklist. They are historical controls for comparing
unlike sources.

\subsection{Test one: who is the recipient represented by the source?}

Recipient identity is the fastest way to expose false continuity. Mark's
restored man, the child in the disputed textual history of Mark~9, a
catechumen renouncing Satan, a person described by a medieval rubric as
\emph{energumenus}, a candidate admitted to an ecclesiastical grade, a church
or object named in a blessing, and the person contemplated by canon~1172 are
not one documentary subject.

The biblical narratives represent persons within plots of proclamation,
conflict, restoration, discipleship, and the arrival of God's reign. They do
not classify a population through a later canonical category. Initiatory
sources address candidates for incorporation into the Church. Their
exorcistic language belongs to conversion and Baptism and does not imply that
every catechumen is being diagnosed as possessed. Ordination witnesses address
candidates for ministry rather than afflicted recipients. Medieval
sacramentaries and pontificals may place several recipient classes in one
book, but book proximity does not erase the rubrical difference.

The 1614 Roman unit and its successors provide a more bounded documentary
setting for a person represented as afflicted in the manner contemplated by
that book. Current Latin law is narrower still in juridical function: canon
1172 governs exorcisms \emph{in obsessos}, and the applicable ritual and
territorial implementation supply further controls. Neither the canon nor the
book retrospectively renames every ancient or medieval recipient.

This test yields a positive continuity and a decisive discontinuity. Across
many sources, a human being is represented as needing divine help and
ecclesial care. The source-specific identity of that person changes
substantially. Human dignity can continue without one unchanged recipient
category.

\subsection{Test two: what act does the source represent?}

The word family ``exorcism'' can name or neighbor very different acts:
Christ's command in a Gospel narrative, an apologist's public challenge,
renunciation before Baptism, prayer over a catechumen, an ordination act, a
blessing, prayer over an afflicted person, a legal reservation, or a modern
liturgical sacramental. Similar vocabulary is evidence for comparison, not a
permission to merge.

Genre supplies the first control. Narrative tells what a narrator represents
as occurring and why it matters in the plot. Apology argues before an
audience and may use public confidence as part of persuasion. A church order
or sacramentary supplies normative or reconstructed liturgical material. An
ordinal represents admission to a grade. A canon regulates conduct. A ritual
book supplies an authorized liturgical form. A clinical manual classifies
presentations for diagnosis and care within medicine. A safeguarding source
names harms and duties within a territorial civil or professional setting.

The historical continuity lies at a higher level than one procedure:
Christians pray to God, confess Christ's victory, reject magic, and order
public ministry within the Church. The acts by which those convictions are
expressed are not invariant. This distinction explains why the present study
can be comprehensive without reproducing formulas. Formula recovery would
deepen a textual history; it would not by itself settle identity, authority,
use, efficacy, or lawful present action.

\subsection{Test three: who is the minister, and what makes the act theirs?}

Ministerial vocabulary is especially vulnerable to anachronism. The Gospels
represent Jesus' authority and delegated mission. Apologists sometimes stress
ordinary Christians rather than a reserved professional class. Eusebius'
quotation of Cornelius names exorcists among Roman personnel but gives a
combined numerical total and does not define their duties. Laodicea~26
restricts unauthorized persons in its own disciplinary setting. The
\emph{Apostolic Constitutions} distinguishes charismatic action from an
ordination conferred for that purpose. Western ordinals represent admission
to the exorcist grade; initiatory books assign actions according to their own
ritual order.

Those pieces cannot be arranged as a simple career ladder ending in the
modern diocesan exorcist. The 1972 reform itself proves that office history
changed: \emph{Ministeria quaedam} suppressed the former minor-order structure
as such, retained lector and acolyte as universal instituted ministries, and
treated any additional regional ministry as a separate conference request to
the Apostolic See. The 2021 amendment of canon~230 §1 changed eligibility for
the named instituted ministries; it did not create a universal ministry of
exorcist.

Canon~1172 supplies the governing Latin rule for the act it names: no one may
lawfully pronounce exorcisms over the possessed without the local ordinary's
particular and express permission, granted only to a priest possessing the
stated qualities. This is a permission rule, not a historical definition of
every exorcist and not proof that every diocese constitutes a permanent
office. U.S. implementation may describe appointment by a diocesan bishop on
a stable or case-specific basis; that is a territorial implementation and
must not be substituted silently for the exact universal wording.

The durable thread is accountable ministry rather than a stable job
description. Claims to spiritual power do not self-authenticate. The
institutional form of accountability, the minister named, and the act reserved
change across sources.

\subsection{Test four: what kind of book or artifact carries the evidence?}

Books create powerful optical illusions. A carefully copied manuscript can
look universally normative. A printed typical edition can look universally
used. A title repeated across editions can look textually motionless. A modern
translation can look like the universal Latin act. Artifact criticism breaks
those illusions.

The Gellone and Kraków manuscripts show placement and differentiated units in
named books. They do not establish frequency, geographical saturation, or one
archetypal practice. Vogel--Elze's PRG edition and the Vendôme manuscript
permit a bounded edition-to-manuscript comparison; they do not turn every
pontifical into a witness to the same local use. Parker MS~79 provides a
late-medieval English ordinal comparison, not representative statistics for
England or the West.

The Roman edition sequence demonstrates real continuity and change at a
different scale. Direct page controls establish the 1614 unit through its
terminal page; the next page's title remains independently located but was not
recovered as a usable exact image in the present 29 July 2026 recheck. The
1872 state remains continuous in architecture; the 1925 edition visibly
divides title XI into chapters; the 1952 typical-edition location is known at
lower textual depth; official acts distinguish the 1998 decree, 1999 physical
issue, and 2004 emended reprint. The history is neither one unchanged book nor
a blank before the latest edition.

Each artifact answers only the question its bytes and identity can bear.
Title, page boundary, chapter architecture, promulgating act, translation
approval, confirmation, and territorial implementation are separate facts.
They should remain separate even when a reader eventually encounters them in
one bound volume.

\subsection{Test five: what authority does the source possess?}

Authority is not a ladder on which every later source simply outranks every
earlier one. It is competence related to a question. Scripture normatively
proclaims the canonical witness but does not contain the text of canon~1172.
A father or medieval theologian may be a major reception witness without
legislating current discipline. A typical liturgical book regulates its act
without supplying clinical diagnosis. The Latin Code governs the Latin Church
and does not become Eastern common law. A conference page establishes
identified territorial implementation and guidance, not universal law. WHO
classification governs clinical description within professional scope and
does not adjudicate a preternatural cause. Civil safeguarding guidance
identifies harms and duties in its jurisdiction without settling Catholic
demonology.

This authority test also exposes false appeals to ``the Church.'' The
Catechism, a dicastery decree, a Code canon, a conference FAQ, an exorcist's
memoir, and a diocesan practice are not interchangeable speakers. A claim
must name the authority, jurisdiction, date, and kind of act. The current Latin
rule remains the 1983 Code as amended through the stated as-of date; the
official pages checked on 29 July 2026 retain canons~1 and 1172 in the form
used here. The official published authentic-interpretations register checked
for this study still supplies no entry for canon~1172. That negative register
check is bounded and does not replace fact-specific legal research.

Historical age does not create present authority, and present authority does
not alter historical evidence. The 1614 book's age cannot authorize private
use; canon~1172 cannot be projected backward as the legal grammar of the
Gospels or outward as the law of an Eastern Catholic Church.

\subsection{Test six: what would count as evidence for practice or efficacy?}

Normative evidence is abundant because institutions preserve books and acts.
Practice evidence is thinner and more vulnerable. A rubric establishes an
available or prescribed form. An ordinal establishes a represented admission.
A canon establishes regulation. None counts actual events or outcomes.
Ownership does not prove use; use does not prove authorization; authorization
does not prove the alleged cause; a narrated outcome does not prove efficacy.

The Loudun module shows why named cases require dependency analysis. Royal
process, defense writing, polemical pamphlet, ecclesiastical voice, later
edition, and cultural retelling may concern the same events while performing
different argumentative work. Their multiplication is not automatic
corroboration. The controlled dossier can establish acts, allegations,
institutional conflict, and reception. It cannot establish possession,
diagnosis, efficacy, guilt beyond the bounded legal record, or the truth of
every narrated incident.

A representative history of practice would require archival series rather
than a chain of famous stories: diocesan faculties and correspondence,
inventories, visitation or seminary records, dated case files with lawful
access, and evidence about preservation bias. Comparative efficacy would
require defined outcomes, denominators, follow-up, and controls no ritual book
or memoir supplies. The present study therefore refuses prevalence and success
rates not because the questions are unimportant, but because the controlled
corpus is not designed to answer them.

The same discipline applies to clinical evidence. The WHO 2024 CDDR describes
possession-trance disorder and relevant differentials within a diagnostic
framework. It neither supplies a ritual criterion nor proves or disproves a
spiritual cause. A clinical workup can identify, treat, or exclude conditions
within medicine; an unexplained remainder is not thereby converted into
positive evidence of possession.

\subsection{Test seven: what form of care remains obligatory under uncertainty?}

The final test is historical, theological, and moral. Sources differ sharply
about recipient, act, minister, and authority, but uncertainty never converts
a person into an experiment. The U.S. conference material checked again on 29
July 2026 continues to emphasize authorization, professional evaluation,
consent where possible, confidentiality, cultural attention, and avoidance of
isolated ministry. Those statements are U.S. guidance and implementation, not
one universal civil or clinical protocol.

The WHO control retains the clinical boundary used in this study. The 2026
Crown Prosecution Service guidance remains a distinct England-and-Wales
prosecution and safeguarding control: it recognizes spiritual or ritualistic
abuse linked to accusations involving witchcraft, spirit possession, or
demonic influence. It does not establish Catholic prevalence. It does
establish that religious interpretation can become a vehicle of coercion or
harm and that good intention is not a safeguarding defense.

Care therefore does not wait for metaphysical agreement. Immediate danger
requires the appropriate emergency or protective response. Medical,
neurological, psychiatric, sleep, substance, trauma, and cultural
differentials remain within competent professional assessment. Prescribed
treatment should not be stopped on a minister's instruction. Consent is
ongoing; a person with capacity may refuse or withdraw from a non-emergency
religious inquiry or rite, and refusal must not be renamed spiritual
resistance. Privacy does not excuse concealment of abuse or override lawful
protective disclosure.

Ordinary Catholic care likewise continues whether or not a reserved rite is
authorized: prayer addressed to God, sacramental life, material assistance,
truthful listening, and accountable pastoral accompaniment. The historical
continuity most worth preserving is therefore not dramatic technique. It is
the subordination of ministry to Christ, truth, ecclesial order, and the
dignity of the person.

\subsection{The comparative result}

\begin{threestable}{Test}{Continuity responsibly affirmed}{Discontinuity or unknown preserved}
Recipient & A person in need of divine help and human care remains a neighbor. & Catechumen, afflicted recipient, office candidate, place, and canon~1172 subject are not one category.\\
Act & Christian prayer confesses God's sovereignty and Christ's victory. & Narrative command, initiation, ordination, blessing, major exorcism, and clinical assessment are different acts.\\
Minister & Public ministry requires ecclesial accountability. & Charism, ancient title, minor order, priesthood, office, and particular permission are not identical.\\
Book & Traditions are transmitted in identifiable artifacts. & Manuscript, critical edition, typical book, translation, and local implementation answer different questions.\\
Authority & Sources possess real but bounded competence. & Scripture, father, theologian, canon, dicastery, conference, clinician, and civil authority do not speak interchangeably.\\
Evidence & Acts, claims, structures, and receptions can be established. & Frequency, cause, efficacy, and representativeness remain unproved without their own evidence.\\
Care & Truth, consent, safety, competent care, and dignity remain obligatory. & Uncertainty never authorizes coercion, spectacle, amateur testing, or treatment displacement.\\
\end{threestable}

The seven tests change the governing thesis from a slogan into a reproducible
historical judgment. Continuity without identity means continuity must be
located at the level supported by the sources, while every changed recipient,
act, minister, book, authority, evidence class, and duty of care remains
visible. That result is neither skeptical evasion nor a pedigree for private
practice. It is the substantive history required before present discipline can
be described honestly.
